\chapter{Low Libido (Female Hypoactive Sexual Desire Disorder)}

\section*{Definition}
Persistent or recurrent absence or deficiency of sexual fantasies, thoughts, and desire for sexual activity, causing marked distress or interpersonal difficulty.

\section*{What Happens in Your Body}
Female sexual desire involves a complex interplay of neurological, hormonal, and psychosocial factors. Testosterone (produced by ovaries and adrenals) is the primary hormone driving sexual desire in women; its levels decline with age and fall significantly after menopause. Oestrogen maintains vaginal health and lubrication; its decline causes dyspareunia, which secondarily reduces desire. Dopamine (reward/anticipation) and serotonin (inhibition) balance sexual motivation. Menopausal hormonal shifts tilt this balance toward inhibition, while sleep deprivation, fatigue, and mood changes further suppress desire.

\section*{Causes}
\begin{itemize}
  \item Menopausal testosterone and oestrogen decline
  \item Psychosocial factors (relationship conflict, stress, body image)
  \item Depression and anxiety
  \item Medication side effects (SSRIs, SNRIs, antipsychotics, hormonal contraceptives)
  \item Vaginal dryness and dyspareunia (fear of pain)
  \item Fatigue and sleep disturbance
  \item Hypothyroidism
  \item Hyperprolactinaemia
  \item Bilateral oophorectomy (surgical menopause)
  \item Chronic illness (diabetes, cardiovascular disease, cancer)
\end{itemize}

\section*{When to See a Doctor}
See your doctor if:
\begin{itemize}
  \item Symptoms are severe or getting worse over time
  \item Low libido causing personal distress, affecting relationships, or associated with depression, pain during sex, or other menopausal symptoms
  \item You have other concerning symptoms such as fever, weight loss, or bleeding
\end{itemize}

\section*{Related Symptoms}
\begin{itemize}
  \item Vaginal dryness
  \item Painful sex
  \item Fatigue
  \item Depression
  \item Anxiety
\end{itemize}

\textbf{Important:} If this symptom persists, your doctor will check for serious underlying causes.

\section*{Self-Care / Home Management}
\begin{itemize}
  \item Rest and avoid activities that make it worse.
  \item Maintain adequate hydration and a balanced diet.
  \item Over-the-counter remedies may provide symptomatic relief (consult a pharmacist or doctor).
  \item Monitor symptoms and seek professional advice if they persist or worsen.
  \item Avoid self-medicating with prescription medications without medical guidance.
\end{itemize}

\section*{How Your Doctor Will Evaluate You}
\begin{itemize}
  \item A thorough history and physical examination are the first steps in evaluation.
  \item Laboratory tests (blood work, urinalysis) may be ordered based on clinical suspicion.
  \item Imaging studies (X-ray, ultrasound, CT, MRI) may be indicated when structural pathology is suspected.
  \item Specialist referral is arranged when the diagnosis remains uncertain or requires specific expertise.
\end{itemize}

\section*{Treatment Options}
\begin{itemize}
  \item Treatment targets the underlying cause identified through diagnostic evaluation.
  \item Supportive care and symptom management are provided while awaiting definitive therapy.
  \item Medications, lifestyle modifications, and physical therapies may be prescribed as appropriate.
  \item Follow-up ensures treatment effectiveness and allows timely adjustments to the care plan.
\end{itemize}

\clearpage